\chapter{الزخم الزاوي، واللف، والإحصاء}

\section{من الدوران العادي إلى اللف الذاتي}

الزخم الزاوي هو الموضوع من الناحية الشكلية، لكن اللف هو المقصود الحقيقي. واللف إحدى أبسط خصائص الجسيمات الأولية وأكثرها منافاة للحدس، وتنبني بنيته على رياضيات الزخم الزاوي.

الزخم الزاوي متجه. واتجاهه يقع على امتداد محور الدوران. وعلى طول ذلك المحور اتجاهان ممكنان، لذا نحتاج إلى اصطلاح للاختيار بينهما: اثنِ أصابع اليد اليمنى في جهة الدوران، فيشير الإبهام إلى اتجاه متجه الزخم الزاوي. وهذا تعريف لا مبرهنة.

يعتمد المقدار على كمية الكتلة الموجودة، وعلى كيفية توزع تلك الكتلة، وعلى حجم النظام، وعلى سرعته الزاوية. ويُجمع توزع الكتلة والحجم في عزم القصور الذاتي. ولذا يمكن لجسم كبير وجسم صغير أن يحملا الزخم الزاوي نفسه وهما يدوران بمعدلات مختلفة جداً وبطاقات مختلفة جداً.

يمكن لجسم مركب عادي أن يحمل نوعين مختلفين مفهوماً من الزخم الزاوي. فالكوب الذي يتحرك في دائرة حول محور خارجي يملك زخماً زاويّاً لأن مركز كتلته يتحرك حول ذلك المحور، حتى لو لم يكن الكوب يدور حول أي محور يمر عبره. وهذا هو الزخم الزاوي المداري. أما اللف فيعني الدوران حول محور داخلي.

ويستخدم التعريف الأدق إطار مركز الكتلة. انتقل إلى الإطار الذي يكون فيه مركز الكتلة ساكناً، بحيث ينعدم الزخم الخطي العادي الكلي. وأي زخم زاوي يبقى في ذلك الإطار هو زخم اللف. ولا يلزم أن يكون الجسم ساكناً داخلياً؛ بل يكفي أن يكون مركز كتلته ساكناً.

في حالة كرة السلة لا يسبب هذا التمييز أي صعوبة. فالتعرّف على كرة سلة ساكنة وعلى كرة السلة نفسها بعد جعلها تدور بوصفهما الجسم نفسه في حالتين مختلفتين أمر واضح. غير أن ميكانيكا الكم تعقّد هذه الصورة لأن الزخم الزاوي متكمِّم. فلا توجد متتالية متصلة من قيم الزخم الزاوي المسموح بها تصل حالةً ذات زخم زاوي صفري بأول حالة دورانية. ومن ثم يمكن للمرء أن يسأل: هل أول حالة دورانية هي الجسم نفسه في حركة، أم جسم مختلف تماماً؟

الصياغة اصطلاحية جزئياً، لكن السؤال الفيزيائي دقيق: ما مقدار الطاقة المطلوبة لبلوغ أول حالة دورانية غير صفرية؟ إذا كانت الفجوة صغيرة، ظل الجسم المثار قابلاً للتعرّف عليه بسهولة. وإذا كانت الفجوة هائلة، فقد تغيّر الإثارة الجسم جذرياً. فمحاولة إعطاء ذرة إثارة دورانية أكبر مما ينبغي قد تقذف إلكتروناتها وتدمّر الذرة، بدلاً من أن تترك الذرة نفسها تدور بهدوء.

والآن لننظر في إلكترون، مع تنحية لفه الذاتي المعروف جانباً مؤقتاً. هل يمكن جعل إلكترون يدور حول محور داخلي؟ النقطة الرياضية لا تملك بنية داخلية مرئية كي تدور، لكن التجربة لا تثبت حجم الإلكترون بأن تزوّدنا بسطح صلب صغير يمكن مشاهدته وهو يدور. فإذا كان للإلكترون حجم متناهٍ صغير جداً، فقد تتطلب أول إثارة دورانية إضافية له طاقةً كبيرة جداً بحيث يتعذر بلوغها. ولأن الطاقة ترتفع كلما صغر الحجم، فقد تكون الطاقة المطلوبة فلكية، وربما حتى من رتبة طاقة بلانك. وهذا التقدير مبدئيٌّ عن قصد.\footnote{مقطع قصير يلي هذا التقدير مباشرة ليس مفهوماً بما يكفي ليدعم إعادة بناء مضمونية أخرى.}

وبناءً على ذلك، يُعامَل مقدار اللف في الجسيم الأولي على أنه سمة ثابتة. فالإلكترونات لا تأتي على هيئة مدى من مقادير اللف الذاتي المتزايدة تباعاً. وإذا حملت إثارة تشبه الإلكترون من سائر الوجوه لفاً كلياً مختلفاً، فستُصنَّف جسيمًا مختلفًا.

المقدار الثابت لا يثبت الاتجاه. تخيّل جسيمًا يُسمّى \emph{سبينترون}، ويُعرَّف جزئياً بمقدار زخمه الزاوي. فإذا كان لفه يشير على طول محور ما، فإن تناظر الدوران يقتضي أن النوع نفسه من الجسيم يمكن أيضاً أن يكون لفه متجهاً على طول محور آخر. ولا يوجد أي اتجاه مكاني مفضَّل جوهرياً. يبقى مقدار اللف ثابتاً، بينما يمكن لاتجاهه أن يتغير. وهكذا تتضمن المسألة عنصرين يبدوان متعارضين: مقداراً متكمَّماً واتجاهات قابلة للدوران بحرية. وستُظهر جبرية الزخم الزاوي كيف يتعايشان.

\begin{classroomqa}
\lecturetimestamp{00:12:25}
\audiencequestion{لماذا يتطلب تدوير جسيم صغير طاقة؟}
\lecturerresponse{لجسم كتلته \(M\) ونصف قطره المميِّز \(R\)، يكون عزم القصور الذاتي من رتبة
\begin{equation}
I\sim cMR^2,
\end{equation}
حيث يعتمد المعامل العددي \(c\) على الشكل وعلى كيفية توزع الكتلة. وإذا تحركت المادة القريبة من الخارج بسرعة \(v\)، فإن زخمها العادي يكون من رتبة \(Mv\)، والزخم الزاوي هو الزخم مضروباً في المسافة:
\begin{equation}
L\sim MvR.
\end{equation}
وبصورة مكافئة، \(L=I\omega\) مع \(v\sim\omega R\). وطاقة الحركة الدورانية هي
\begin{equation}
E_{\mathrm{rot}}=\frac{L^2}{2I}.
\end{equation}
وعند ثبات \(M\) وثبات \(L\)،
\begin{equation}
E_{\mathrm{rot}}\sim\frac{L^2}{2cMR^2}\propto\frac{1}{R^2}.
\end{equation}
فكلما صغر الجسم زادت الطاقة التي يكلّفها مقدار معيّن من الزخم الزاوي. وهذا الاعتماد على معكوس المربع كلاسيكي. أما الحقيقة الكمومية الإضافية فهي أن الزخم الزاوي يتغير في خطوات منفصلة، لذا لا يمكن جعل أول خطوة مسموح بها صغيرة على نحو اعتباطي.}
\end{classroomqa}

الرياضيات التي سنستخدمها الآن مجردة، وفي البداية بعيدة نوعاً ما عن الحدس. وتنبع أهميتها من أن حقائق قابلة للاختبار تجريبياً تنشأ من مجموعة صغيرة من القواعد الجبرية.

\section{الزخم الزاوي المداري}

لنبدأ بجسيم واحد يتحرك في مدار دائري حول مبدأ الإحداثيات. مقدار زخمه الزاوي يساوي الزخم مضروباً في البعد العمودي عن المبدأ. وفي مستوى يمكن اعتبار أحد اتجاهي الدوران موجباً والاتجاه المعاكس سالباً. وتسجل الإشارة اتجاه متجه الزخم الزاوي بالنسبة إلى المحور المختار.

هذا ليس لفاً بعد. فللجسيم زخم زاوي لأنه يملك زخماً عادياً. ضع الآن جسيمين عند طرفين متقابلين لقطر، ودعهما يتحركان في جهة الدوران نفسها. عندئذ يتلاشى زخماهما العاديان، فيكون مركز الكتلة ساكناً، لكن زخميهما الزاويين يشيران في الاتجاه نفسه ويُضاف أحدهما إلى الآخر. للنظام زخم كلي صفري وزخم زاوي كلي غير صفري: نموذج كلاسيكي للَّف.

ليكن موضع جسيم بالنسبة إلى المبدأ
\begin{equation}
\vec r=(x,y,z)=(x_1,x_2,x_3),
\end{equation}
وليكن زخمه
\begin{equation}
\vec p=(p_x,p_y,p_z)=(p_1,p_2,p_3).
\end{equation}
فالزخم الزاوي هو حاصل الضرب الاتجاهي
\begin{equation}
\vec L=\vec r\times\vec p.
\end{equation}
وعكس العاملين يقلب المتجه: \(\vec p\times\vec r=-\vec r\times\vec p\). والاصطلاح المعياري هو \(\vec r\times\vec p\)، ويُحدَّد اتجاهه بقاعدة اليد اليمنى.\footnote{\raggedright إطار من فيديو المحاضرة \nolinkurl{lecture_07_figure_02.png}، 00:18:49.\par}

\begin{figure}[h]
\centering
\includegraphics[width=0.72\textwidth]{lecture_07_figure_02.png}
\par
\begin{tikzpicture}[scale=1.0]
\draw[->] (-0.6,0) -- (4.8,0) node[right] {$x$};
\draw[->] (0,-0.6) -- (0,3.8) node[above] {$y$};
\coordinate (P) at (3.2,2.3);
\draw[->,thick] (0,0) -- (P) node[midway,above left] {$\vec r$};
\draw[dashed] (P) -- (3.2,0);
\draw[dashed] (P) -- (0,2.3);
\draw[->,thick] (P) -- ++(-1.0,1.25) node[above] {$\vec p$};
\node at (4.15,2.75) {$\vec L=\vec r\times\vec p$};
\end{tikzpicture}
\caption{بالنسبة إلى جسيم عند الموضع \(\vec r\) وله زخم \(\vec p\)، يكون الزخم الزاوي \(\vec L=\vec r\times\vec p\) عمودياً على المستوى الذي يحددهما.}
\lectureframe{00:18:49}
\end{figure}

في المركبات الديكارتية،
\begin{equation}
\begin{aligned}
L_x&=yp_z-zp_y,\qquad
L_y=zp_x-xp_z,\\
L_z&=xp_y-yp_x.
\end{aligned}
\end{equation}
ويكفي أن تتذكر الصيغة الأولى وأن تدوّر الرموز
\[
x\longrightarrow y\longrightarrow z\longrightarrow x.
\]
فبتدوير تسمية المركبة والتسميات الواقعة في الطرف الأيمن معاً، تحصل على الصيغة التالية من دون تغيير الإشارة النسبية. وهذه القاعدة التذكيرية مفيدة تحديداً لأن من السهل عكس الإشارة من غير قصد.

بالنسبة إلى جسيم نقطي، \(\vec r\times\vec p\) هو الزخم الزاوي المداري حول المبدأ المختار. أما في نظام مركب، فنجمع متجهات الزخم الزاوي لجميع المكونات:
\begin{equation}
\vec L_{\mathrm{tot}}=\sum_a \vec r_a\times\vec p_a.
\end{equation}
وعندما يُحسَب هذا المجموع في إطار مركز الكتلة، يمكن أن يبقى غير صفري رغم أن
\[
\sum_a\vec p_a=0.
\]
ذلك الزخم الزاوي المتبقي هو لف النظام المركب.

\section{المؤثرات الكمومية والمبدلات القانونية}

في ميكانيكا الكم، تُمثَّل الكميات القابلة للرصد بمؤثرات. لذلك فإن الموضع، والزخم، وكل مركبة من مركبات الزخم الزاوي هي مؤثرات. ولسنا بحاجة هنا إلى تحقيقها التفصيلي على هيئة مؤثرات تفاضلية. وتحتوي علاقات التبادل على معلومات كافية لاشتقاق طيف الزخم الزاوي.

لمؤثرين \(A\) و\(B\)،
\begin{equation}
[A,B]=AB-BA.
\end{equation}
إذا كان مرصودان يتبادلان، أمكن تحديدهما معاً. وإذا لم يتبادلا، منعت ميكانيكا الكم تحديدهما الدقيق معاً.

يمكن قياس إحداثيات الموضع الثلاثة معاً، لذا تتبادل مركبات الموضع:
\begin{equation}
[x_i,x_j]=0.
\end{equation}
وعلى وجه الخصوص، كل مؤثر يتبادل مع نفسه. ولا معنى لأن يكون مرصود غير متوافق مع قياس آخر لذلك المرصود نفسه.

وبالمثل، تتبادل مركبات الزخم:
\begin{equation}
[p_i,p_j]=0.
\end{equation}
وتتبادل مركبة من الموضع مع مركبة من الزخم في اتجاه مختلف، لكنها لا تتبادل مع مركبة الزخم المناظرة لها. ويمكن كتابة العلاقات القانونية الثلاث معاً على الصورة
\begin{equation}
[x_i,x_j]=0,\qquad
[p_i,p_j]=0,\qquad
[x_i,p_j]=i\hbar\,\delta_{ij}.
\end{equation}
وهكذا، على سبيل المثال،
\[
[x,p_y]=0,\qquad [x,p_x]=i\hbar,
\]
وعكس الترتيب يعكس الإشارة:
\begin{equation}
[p_j,x_i]=-i\hbar\,\delta_{ij}.
\end{equation}

للعزم الزاوي وثابت بلانك الأبعاد نفسها. فالعزم الزاوي هو المسافة مضروبة في الزخم، بينما تأخذ علاقة عدم اليقين الصورة البعدية
\begin{equation}
\Delta x\,\Delta p\sim\hbar
\end{equation}
باستثناء العامل العددي المألوف في المتباينة الدقيقة. لذلك فمن الطبيعي أن يكون العزم الزاوي مكمّى بوحدات من \(\hbar\)، مع أن هذا التكميم لم يُشتق بعد.

\section{اشتقاق جبر العزم الزاوي}

المهمة الأولى هي حساب مبدّلة مركبتين من العزم الزاوي. انطلاقًا من
\[
L_x=yp_z-zp_y,\qquad
L_y=zp_x-xp_z,
\]
نحصل على
\begin{align}
[L_x,L_y]
&=[yp_z-zp_y,\;zp_x-xp_z] \\
&=[yp_z,zp_x]-[yp_z,xp_z] \notag\\
&\quad{}-[zp_y,zp_x]+[zp_y,xp_z].
\end{align}

\begin{classroomqa}
\lecturetimestamp{00:31:41}
\audiencequestion{هل يمكن تذكّر صيغ المركبات بتدوير الرموز من اليسار إلى اليمين بدلًا من تدوير المعادلات الثلاث من الأعلى إلى الأسفل؟}
\lecturerresponse{نعم. أيٌّ من طريقتي التذكّر القائمتين على التدوير مقبول. والغرض الوحيد منها هو توليد المركبات بالترتيب الدوري الصحيح مع تتبّع الإشارة.}
\end{classroomqa}

يمكن فهم المبدّلة بالسؤال عمّا إذا كان يمكن تمرير العوامل في أحد الجداءين عبر العوامل في الآخر. في
\[
[yp_z,xp_z],
\]
يتبادل المؤثر \(y\) مع كلٍّ من \(x\) و\(p_z\)، بينما يتبادل \(p_z\) مع كلٍّ من \(x\) ونفسه. لا يوجد عائق، لذا ينعدم هذا الحد. ويعطي الاستدلال نفسه
\[
[zp_y,zp_x]=0.
\]
ولا يبقى إلا زوجان من الجداءات.

أما في الزوج الأول، فالعائق الوحيد هو مرور \(p_z\) عبر \(z\):
\begin{align}
[yp_z,zp_x]
&=yp_zzp_x-zp_xyp_z \\
&=y(p_zz-zp_z)p_x \\
&=y[p_z,z]p_x \\
&=-i\hbar\,yp_x.
\end{align}
وترتيب المبدّلة القانونية مهم:
\[
[p_z,z]=-[z,p_z]=-i\hbar.
\]

أما الزوج الباقي الآخر، فإن إشارتي السالب في تعريفَي \(L_x\) و\(L_y\) تعطيان عند ضربهما إشارة موجبة. والعائق الوحيد هو مرور \(z\) عبر \(p_z\):
\begin{align}
[zp_y,xp_z]
&=zp_yxp_z-xp_zzp_y \\
&=x(zp_z-p_zz)p_y \\
&=x[z,p_z]p_y \\
&=i\hbar\,xp_y.
\end{align}
وبجمع المساهمتين،
\begin{align}
[L_x,L_y]
&=-i\hbar\,yp_x+i\hbar\,xp_y \\
&=i\hbar(xp_y-yp_x) \\
&=i\hbar L_z.
\end{align}

هذا هو التبسيط المهم. إن أخذ مبدّلة بين مركبتين من العزم الزاوي يُنتج المركبة الثالثة؛ ولا يبقى أي مؤثر موضع أو زخم. ويعطي التبديل الدوري
\begin{equation}
\begin{aligned}
[L_x,L_y]&=i\hbar L_z,\qquad
[L_y,L_z]=i\hbar L_x,\\
[L_z,L_x]&=i\hbar L_y.
\end{aligned}
\end{equation}
وبصورة مكافئة،
\begin{equation}
[L_i,L_j]=i\hbar\,\epsilon_{ijk}L_k.
\end{equation}

ينغلق هذا الجبر على مؤثرات العزم الزاوي نفسها. وبمجرد الحصول على هذه العلاقات، يمكن تنحية \(x_i\) و\(p_i\) اللذين استُخدما لاشتقاقها جانبًا. وينطبق الجبر نفسه عندما يكون العزم الزاوي الكلي مجموعًا على عدد كبير من المكوّنات. وعلى وجه الخصوص، فإنه ينطبق في إطار مركز الكتلة، حيث يكون العزم الزاوي المتبقي هو اللفّ المغزلي. ولذلك يمكن استخدام مبدّلات العزم الزاوي على نحو تجريدي من دون أن نبني أولًا اللفّ المغزلي بوصفه حركة أجزاء مكوِّنة صغيرة.

اختيار المحاور \(x\) و\(y\) و\(z\) اعتباطي. فإذا أجرينا دورانًا إلى أي مجموعة أخرى من المحاور الديكارتية القصورية المتعامدة فيما بينها، فإن المركبات ذات الشرطة تحقق العلاقات نفسها:
\begin{equation}
[L'_x,L'_y]=i\hbar L'_z,
\end{equation}
مع نظائرها الدورية. هذا الجبر ثابت تحت الدوران، لذا لا يمكن أن تعتمد التنبؤات الفيزيائية على المحاور القصورية التي اختيرت للحساب.

\begin{classroomqa}
\lecturetimestamp{00:39:09}
\audiencequestion{هل العزم الزاوي المغزلي مستقل عن انتقالات محاور الإحداثيات وكذلك عن دوراناتها؟}
\lecturerresponse{نعم. اللف الذاتي مستقل عن الانتقال. ويمكن أن يتغير العزم الزاوي المداري بالنسبة إلى مبدأ ما عند نقل ذلك المبدأ، لكن اللف الباقي في إطار مركز الكتلة لا يعتمد على ذلك الانتقال.}
\end{classroomqa}

يمكن للجبر المجرّد أن يولّد عددًا لا ينتهي من المتطابقات، ومعظمها غير ذي أهمية فيزيائية. وتظهر قيمته عندما تصبح متطابقة ما حقيقة قابلة للقياس أو تكشف شيئًا عن المعنى الإجرائي للقياس. وحتى هذه النقطة، لا يزال البناء في معظمه جبريًا. وتستخلص الخطوة التالية البنية المتقطعة للَّف.

\begin{classroomqa}
\lecturetimestamp{00:40:42}
\audiencequestion{ماذا لو دارت المحاور المختارة \(x\) و\(y\) و\(z\) مع الجسيمات؟}
\lecturerresponse{يقتصر النقاش على الأطر العطالية. فالنظام الإحداثي الدوار غير عطالي ويُدخل تأثيرات مثل قوتي الطرد المركزي وكوريوليس. ولا يبقى العزم الزاوي ثابتًا عند الانتقال إلى إطار مرافق للدوران: فالنظام الذي يبدو دوارًا بوضوح في إطار عطالي قد يبدو ساكنًا في إطار يدور معه. أمّا الثبات الدوراني المستخدم هنا فيتعلق بمحاور ديكارتية عطالية مختلفة التوجّه، لا بمحاور تدور مع الزمن.}
\end{classroomqa}

\section{مؤثرا الرفع والخفض ومحور التكميم}

عرّف التركيبين التاليين
\begin{equation}
L_+=L_x+iL_y,
\qquad
L_-=L_x-iL_y.
\end{equation}
وهما نظيران لمؤثري الخلق والإفناء في المذبذب التوافقي. ويتمثل دورهما في رفع وخفض مركبة واحدة مختارة من العزم الزاوي.

لا يوجد ما يميز المحور \(z\)، ولكن لا بد من اختيار محور إذا أريد تعيين مركبة. خذ \(z\) محورًا للتكميم واكتب
\begin{equation}
L_z\lvert m\rangle=m\hbar\,\lvert m\rangle.
\end{equation}
وعليه فإن \(m\) هو قيمة المركبة المقاسة على المحور \(z\) بوحدات \(\hbar\). وغالبًا ما يكون من الملائم أثناء المعالجات الجبرية أن نضع \(\hbar=1\)؛ لكننا سنبقيه هنا صريحًا.

المبدّل \([L_+,L_-]\) يتناسب مع \(L_z\)، لكنه غير لازم للاستدلال المباشر. وما يهم أولًا هو مبدّل كل من مؤثري الرفع والخفض مع \(L_z\). ومن
\[
[L_x,L_z]=-i\hbar L_y,
\qquad
[L_y,L_z]=i\hbar L_x,
\]
نحصل على
\begin{align}
[L_+,L_z]
&=[L_x+iL_y,L_z] \\
&=[L_x,L_z]+i[L_y,L_z] \\
&=-i\hbar L_y+i(i\hbar L_x) \\
&=-\hbar(L_x+iL_y) \\
&=-\hbar L_+.
\end{align}
وبتغيير إشارة الحد \(iL_y\) نحصل على
\begin{equation}
[L_-,L_z]=+\hbar L_-.
\end{equation}
وبصورة مكافئة، إذا عُكس ترتيب كل مبدّل، فإن
\begin{equation}
[L_z,L_+]=\hbar L_+,
\qquad
[L_z,L_-]=-\hbar L_-.
\end{equation}

لا يمكن أن تكون محددة بدقة في وقت واحد إلا مركبة واحدة من العزم الزاوي. فأي مركبتين مختلفتين لا تتبادلان، ولذلك لا يمكن تعيينهما معًا. والمركبة المختارة المحددة بدقة تعرّف محور التكميم؛ واختيار \(z\) هو من باب الملاءمة، لا تفضيلًا فيزيائيًا.

طبّق الآن \(L_+\) على حالة ذاتية لـ \(L_z\). ويمكن إعادة كتابة المبدّل على الصورة
\begin{equation}
L_zL_+=L_+L_z+\hbar L_+.
\end{equation}
ولذلك
\begin{align}
L_z\bigl(L_+\lvert m\rangle\bigr)
&=\bigl(L_+L_z+\hbar L_+\bigr)\lvert m\rangle \\
&=m\hbar L_+\lvert m\rangle+\hbar L_+\lvert m\rangle \\
&=(m+1)\hbar\,L_+\lvert m\rangle.
\end{align}
وشريطة ألا يكون صفريًا، فإن \(L_+\lvert m\rangle\) حالة ذاتية لـ \(L_z\) زادت قيمتها الذاتية بمقدار وحدة واحدة من \(\hbar\). وبالطريقة نفسها،
\begin{equation}
L_z\bigl(L_-\lvert m\rangle\bigr)
=(m-1)\hbar\,L_-\lvert m\rangle.
\end{equation}
وهكذا يخفض \(L_-\) قيمة \(m\) بمقدار واحد.

هذا بالفعل استنتاج غير بديهي. ومهما يكن طيف \(L_z\) في النهاية، فإن القيم المتجاورة فيه تفصلها أعداد صحيحة:
\begin{equation}
\ldots,\ m-2,\ m-1,\ m,\ m+1,\ m+2,\ \ldots
\end{equation}
ولم يُحدَّد الحدّان الطرفيان بعد.

\section{المتعدديات المنتهية ووسم اللف المغزلي}

لنوع معيّن من الجسيمات مقدار ثابت للزخم الزاوي. ويمكن إعادة توجيه زخمه الزاوي، لكن زيادة مقداره بلا حد ستؤدي في النهاية إلى كيان مختلف. لذلك لا بد من وجود أكبر مركبة ممكنة على \(z\) لذلك الجسيم.

كيف يمكن لسلّم ذي مؤثر رفع أن تكون له قمة؟ يوفّر المذبذب التوافقي النموذج. فمؤثر الإفناء فيه يعطي صفراً عندما يبلغ الحالة الدنيا. وهنا يجب أن يعطي مؤثر الرفع صفراً عندما يبلغ الحالة العليا:
\begin{equation}
L_+\lvert m_{\max}\rangle=0.
\end{equation}
إن العبارة القائلة إن \(L_+\) يرفع قيمةً ذاتيةً لا تنطبق إلا إذا كان المتجه الذي ينتجه غير صفري. والحالة ذات الوزن الأعلى هي بالضبط الحالة الاستثنائية التي تنتهي عندها متتالية الرفع.

يحدّد تناظر الدوران الطرف الأدنى. فإذا أمكن للزخم الزاوي أن يصطف بأقصى مقدار على امتداد \(+z\)، أمكن تدوير الكيان الفيزيائي نفسه بحيث يصطف بأقصى مقدار على امتداد \(-z\). ومن ثم
\begin{equation}
m_{\min}=-m_{\max}.
\end{equation}
فالطيف متناظر حول الصفر وتباعده وحدي. وهناك عائلتان ممكنتان.

في العائلة الأولى، يوجد الصفر:
\begin{equation}
\begin{aligned}
m={}&-m_{\max},-m_{\max}+1,\ldots,-1,0,\\
&1,\ldots,m_{\max}-1,m_{\max}.
\end{aligned}
\end{equation}
وعندئذٍ يكون \(m_{\max}\) عددًا صحيحًا، ويكون عدد الحالات هو
\begin{equation}
2m_{\max}+1.
\end{equation}
والقيمة العظمى هي التي تميّز الجسيم وتسمّى لفه المغزلي. وفي الترميز المتعارف عليه،
\begin{equation}
s\equiv m_{\max},
\qquad
\lvert s,m\rangle\equiv\lvert m\rangle
\end{equation}
داخل متعددية ذات لف مغزلي ثابت.\footnote{\raggedright إطار من فيديو المحاضرة \nolinkurl{lecture_07_figure_03.png}، 01:01:08.\par}

\begin{figure}[htbp]
\centering
\includegraphics[width=0.70\textwidth]{lecture_07_figure_03.png}
\caption{ينهي شرط الوزن الأعلى متتالية الرفع، وتشكل قيم \(m\) المسموح بها سلّمًا منتهيًا بتباعد وحدي.}
\lectureframe{01:01:08}
\end{figure}

ليس للجسيم ذي اللف المغزلي الصفري إلا \(m=0\). وللجسيم ذي اللف المغزلي الواحد
\[
m=1,0,-1,
\]
أما الجسيم الذي لفه المغزلي اثنان فله خمس قيم للمركبة. وبوزون هيغز مثال على اللف المغزلي الصفري. وللفوتون وبوزون \(Z\) لف مغزلي مقداره واحد. أما الغرافيتون، إن وُجد، فلفه المغزلي مقداره اثنان. ويمكن للأجسام المركبة أن يكون لها لف مغزلي مقداره ثلاثة وأعلى من ذلك بكثير: فالنوى يمكن أن تفعل ذلك، كما يمكن للأجسام الدوّارة العيانية بما يكفي، مثل كرات السلة والمجرات, أن تحمل زخماً زاويًا بلا حد. وأما ما يغيب عن القائمة المعتادة للجسيمات الأولية فهو جسيم أولي ذو لف مغزلي مقداره ثلاثة.

أما العائلة الأخرى المتناظرة ذات التباعد الوحدي فلا تحتوي على الصفر. وهي تبدأ بـ
\[
m=\frac12,-\frac12
\]
وتمتد بخطوات وحدية:
\begin{equation}
\begin{aligned}
m={}&-m_{\max},-m_{\max}+1,\ldots,-\frac32,-\frac12,\\
&\frac12,\frac32,\ldots,m_{\max}.
\end{aligned}
\end{equation}
وهنا يكون \(m_{\max}\) عددًا نصف صحيح. وأبسط أفرادها له لف مغزلي مقداره \(\tfrac12\)، ثم \(\tfrac32\)، ثم \(\tfrac52\)، وهكذا. ولذلك تأتي متعدديات الزخم الزاوي في عائلتي اللف المغزلي الصحيح واللف المغزلي نصف الصحيح.

\begin{classroomqa}
\lecturetimestamp{01:03:46}
\audiencequestion{هل لمتعددية ذات لف مغزلي نصف صحيح أيضًا قيمة عظمى؟}
\lecturerresponse{نعم. للسلم هنا أيضًا قيمة عليا. فبالنسبة إلى جسيم ذي لف مغزلي \(\tfrac32\) تكون القيمة العظمى هي \(m_{\max}=\tfrac32\)؛ وبالنسبة إلى جسيم ذي لف مغزلي \(\tfrac52\) تكون هي \(m_{\max}=\tfrac52\). والفرق الوحيد هو أن القيمة العظمى وكل قيمة أخرى في المتعددية أعداد نصف صحيحة.}
\end{classroomqa}

\begin{classroomqa}
\lecturetimestamp{01:04:04}
\audiencequestion{ألا توجد جسيمات أولية ذات لف مغزلي ثلاثة أو أعلى؟}
\lecturerresponse{ضمن الإطار الكمي النسبي قيد النقاش، تكون لفوف الجسيمات الأولية المغزلية في القائمة المعتادة هي \(0,\tfrac12,1,\tfrac32,\) و\(2\). أما الأجسام المركبة فليست مقيَّدة بهذا المجال، ويمكن أن يكون لها لف مغزلي ثلاثة أو لف مغزلي أكبر بكثير.}
\end{classroomqa}

\begin{classroomqa}
\lecturetimestamp{01:04:50}
\audiencequestion{هل تمثل المدخلات السالبة على السلم لفًا مغزليًا سالبًا؟}
\lecturerresponse{لا. إن \(m\) السالب هو قيمة للمركبة: فهو يفيد بأن العزم الزاوي يشير في الاتجاه المعاكس بالنسبة إلى المحور المختار. أما لف الجسيم المغزلي فهو القيمة العظمى غير السالبة \(s=m_{\max}\).}
\end{classroomqa}

\begin{classroomqa}
\lecturetimestamp{01:05:10}
\audiencequestion{إذا كانت قيم المركبة متقطعة، فماذا يحدث لكل الاتجاهات الواقعة بينها؟}
\lecturerresponse{ليست هناك جهتان مكانيتان ممكنتان فحسب، ولا ثلاث فحسب. تُمثَّل الاتجاهات البينية بتراكبات كمومية لحالات الأساس \(m\). ويجعل البناء الخاص بلف مغزلي \(\tfrac12\) ذلك صريحًا أدناه.}
\end{classroomqa}

\begin{classroomqa}
\lecturetimestamp{01:05:42}
\audiencequestion{وماذا عن الحركة اللولبية أو الموجات المستقطبة دائريًا؟}
\lecturerresponse{تُبنى الموجة الكهرومغناطيسية المستقطبة دائريًا من فوتونات مستقطبة دائريًا، وتحمل تلك الفوتونات عزمًا زاويًا على طول محور الانتشار. إن لفظة ``لولبي'' تجمع بين الدوران والحركة إلى الأمام، لذا فهي ليست الوصف الصحيح لجسم يُنظر إليه على أنه ساكن. أما بالنسبة إلى جسيم متحرك، فالكمية ذات الصلة هي اللولبية: اتجاه لفه المغزلي بالنسبة إلى اتجاه حركته.}
\end{classroomqa}

\begin{classroomqa}
\lecturetimestamp{01:07:10}
\audiencequestion{هل يوجد أي جسيم أولي معروف ذو لف مغزلي \(\tfrac32\)؟}
\lecturerresponse{لم يُؤكَّد مباشرةً أي جسيم أولي ذي لف مغزلي \(\tfrac32\). وتفترض التناظر الفائق وجود شريك للغرافيتون ذي لف مغزلي \(\tfrac32\) يسمى الغرافيتينو، لكنه لم يُكتشف مباشرةً.}
\end{classroomqa}

\begin{classroomqa}
\lecturetimestamp{01:08:01}
\audiencequestion{هل توقفت مؤثرات العزم الزاوي عن الفعل على فضاء الدالة الموجية أو فضاء الكيت الأصلي؟}
\lecturerresponse{لم يكن قد جرى تحديد أي فضاء لف مغزلي أصلي بعد. وتُستعمل الجبرية لتعيين ذلك الفضاء. وعند تثبيت اللف المغزلي \(s\)، تكون له متجهات أساس موسومة بالقيم المسموح بها
\[
\lvert s,m\rangle,\qquad m=-s,-s+1,\ldots,s.
\]
ويتميّز فضاء اللف المغزلي الداخلي هذا عن الفضاء المعتاد للدوال الموجية المكانية الذي يعمل عليه العزم الزاوي المداري بوصفه \(\mathbf r\times\mathbf p\)، مع أن كليهما يحقق جبرية العزم الزاوي نفسها.}
\end{classroomqa}

\section{مقدار العزم الزاوي}

\begin{classroomqa}
\lecturetimestamp{01:09:32}
\audiencequestion{هل مقدار متجه العزم الزاوي قابل للرصد؟}
\lecturerresponse{نعم. ومن الأنسب استخدام مربعه،
\[
L^2=L_x^2+L_y^2+L_z^2.
\]
وليست قيمته على حالة الوزن الأعلى ببساطة \(s^2\hbar^2\). فحتى عندما تكون \(L_z\) محددة بدقة وعظمى، لا يمكن للمركبتين المستعرضتين غير المتبادلتين أن تكونا كلتاهما صفراً محدداً بدقة، ولذلك تسهمان بمساهمة كمومية إضافية.}
\end{classroomqa}

لحساب تلك المساهمة، اضرب مؤثري السلم بالترتيب الذي يجعل مؤثر الرفع يؤثر أولًا في كيت:
\begin{align}
L_-L_+
&=(L_x-iL_y)(L_x+iL_y) \\
&=L_x^2+L_y^2+iL_xL_y-iL_yL_x \\
&=L_x^2+L_y^2+i[L_x,L_y] \\
&=L_x^2+L_y^2-\hbar L_z.
\end{align}
وتستخدم الخطوة الأخيرة \([L_x,L_y]=i\hbar L_z\). ولذلك
\begin{equation}
L_x^2+L_y^2=L_-L_++\hbar L_z,
\end{equation}
ومن ثم
\begin{equation}
L^2=L_-L_++L_z^2+\hbar L_z.
\end{equation}

طبّق هذه المتطابقة على حالة الوزن الأعلى \(\lvert s,s\rangle=\lvert m_{\max}\rangle\):
\begin{align}
L^2\lvert s,s\rangle
&=\left(L_-L_++L_z^2+\hbar L_z\right)\lvert s,s\rangle \\
&=L_-L_+\lvert s,s\rangle
+s^2\hbar^2\lvert s,s\rangle
+s\hbar^2\lvert s,s\rangle.
\end{align}
يضمحل الحد الأول لأن \(L_+\lvert s,s\rangle=0\). ومن ثم
\begin{equation}
L^2\lvert s,s\rangle
=s(s+1)\hbar^2\lvert s,s\rangle.
\end{equation}
إن الزيادة \(+s\) هي التصحيح الكمومي المرتبط بالتقلبات المستعرضة التي لا مفر منها. وكان متوقعًا كلاسيكيًا \(s^2\hbar^2\) فقط.

تجعل القيم القليلة الأولى هذا التمييز ملموسًا:
\begin{align}
s=0 &: \qquad L^2=0,\\
s=1 &: \qquad L^2=2\hbar^2,\\
s=\frac12 &: \qquad L^2=\frac34\hbar^2.
\end{align}
عند \(s=1000\)، يكون العامل عديم الأبعاد هو \(1000\times1001\)، وهو لا يختلف إلا قليلًا عن \(1000^2\) على مقياس عزم زاوي بهذا الكبر. وبوجه عام،
\begin{equation}
s(s+1)=s^2\left(1+\frac1s\right),
\end{equation}
ولذلك يصبح التصحيح الكمومي النسبي مهمَلًا عند القيم الكبيرة جدًا لـ \(s\). وهكذا تظهر النتيجة الكلاسيكية \(s^2\).

لكل حالة \(m\) في متعددية ذات لف مغزلي ثابت القيمة الذاتية نفسها لـ \(L^2\). إن تغيير \(m\) يغيّر المركبة على طول المحور المختار، لا المقدار الكلي. وهذه هي النسخة الكمومية من تدوير متجه عزم زاوي من دون تغيير طوله.

\begin{classroomqa}
\lecturetimestamp{01:16:43}
\audiencequestion{هل يمكن معرفة \(L^2\) ومركبة \(z\) معًا؟}
\lecturerresponse{نعم. فالمربع الكلي يتبادل مع كل مركبة من مركبات العزم الزاوي:
\[
[L^2,L_i]=0,\qquad i=x,y,z.
\]
وبذلك يمكن تحديد \(L^2\) ومركبة مختارة واحدة، مثل \(L_z\)، معًا، رغم أنه لا يمكن تحديد مركبتين مختلفتين معًا. فلكل إلكترون، على سبيل المثال، \(L^2=\tfrac34\hbar^2\)، بصرف النظر عن اتجاهه.}
\end{classroomqa}

\begin{classroomqa}
\lecturetimestamp{01:17:41}
\audiencequestion{هل يكون \(L^2\) صفراً عند لف مغزلي واحد، أم عند اللف المغزلي الصفري فقط؟}
\lecturerresponse{إن اللف المغزلي الصفري وحده هو الذي له \(L^2=0\). أما اللف المغزلي الواحد فيعطي
\[
L^2=1(1+1)\hbar^2=2\hbar^2,
\]
بينما يعطي اللف المغزلي الثلاثة
\[
L^2=3(3+1)\hbar^2=12\hbar^2.
\]}
\end{classroomqa}

\begin{classroomqa}
\lecturetimestamp{01:18:13}
\audiencequestion{لماذا يكون لكل مستوى \(m\) في المتعددية القيمة الذاتية نفسها لـ \(L^2\)؟}
\lecturerresponse{ابدأ من حالة الوزن الأعلى وأنزلها. وبما أن \(L^2\) يتبادل مع كل مركبة، فهو يتبادل كذلك مع \(L_-\):
\[
[L^2,L_-]=0.
\]
إذا كان
\[
L^2\lvert s,s\rangle=s(s+1)\hbar^2\lvert s,s\rangle,
\]
فحينئذ
\begin{align}
L^2\bigl(L_-\lvert s,s\rangle\bigr)
&=L_-L^2\lvert s,s\rangle \\
&=s(s+1)\hbar^2\,L_-\lvert s,s\rangle.
\end{align}
ومن ثم فإن الحالة التالية نزولًا لها القيمة الذاتية نفسها لـ \(L^2\). ويؤدي تكرار هذه الحجة إلى حمل تلك القيمة الذاتية على طول السلم كله حتى أسفله.}
\end{classroomqa}

ومن ثم فإن قيم \(m\) المتقطعة لا تعني أن مجموعة متقطعة فقط من الاتجاهات المكانية موجودة. فما يزال ثبات الدوران يقتضي حالات تقابل اتجاهات مُدارة. وبالنسبة إلى لف مغزلي \(\tfrac12\)، ستظهر تلك الاتجاهات على هيئة تراكبات لحالتي الأساس \(L_z\) الاثنتين.

\section{حالات اللف المغزلي-\texorpdfstring{\(\tfrac12\)}{1/2} والاتجاه}

إن ثبات الدوران يقتضي أكثر من الحالات الذاتية المتقطعة لـ \(L_z\). إذ يمكن للَّف المغزلي أن يشير على طول أي محور، مع أن قياس مركبته على طول محور مختار لا يعطي إلا القيم المتقطعة المسموح بها. وتُمثَّل الاتجاهات الواقعة بين الحالات الذاتية لمحور \(z\) بتراكبات.

بالنسبة إلى جسيم ذي لف مغزلي \(\tfrac12\)، اقصر الانتباه على درجات حرية اللف المغزلي. فللإلكترون أيضًا حالات مدارية، لذلك فهو بالتأكيد لا يملك حالتين فقط في المجمل؛ إن فضاء لفه المغزلي وحده ثنائي الأبعاد. اختر أساس محور \(z\)
\begin{equation}
L_z\lvert+\rangle=\frac{\hbar}{2}\lvert+\rangle,
\qquad
L_z\lvert-\rangle=-\frac{\hbar}{2}\lvert-\rangle.
\end{equation}
وحالة لف مغزلي عامة هي
\begin{equation}
\lvert\psi\rangle
=
\alpha\lvert+\rangle+\beta\lvert-\rangle,
\end{equation}
حيث إن \(\alpha\) و\(\beta\) عددان مركبان. ويعطي قياس \(L_z\) القيمة \(+\hbar/2\) باحتمال \(|\alpha|^2\) والقيمة \(-\hbar/2\) باحتمال \(|\beta|^2\)، ولذلك يقتضي التطبيع
\begin{equation}
|\alpha|^2+|\beta|^2=1.
\end{equation}

ليس للطور المشترك أي أثر فيزيائي:
\begin{equation}
(\alpha,\beta)\longrightarrow
e^{i\chi}(\alpha,\beta).
\end{equation}
تبدأ السعتان المركبتان بأربعة معلمات حقيقية. يزيل التطبيع واحدًا منها، ويزيل الطور الكلي غير ذي الصلة واحدًا آخر، فيتبقى معلمان حقيقيان فيزيائيان. وهذا هو بالضبط العدد اللازم لتعيين اتجاه بزاويتيه القطبية والسمتية. إن تغيير \(\alpha\) و\(\beta\) يدوّر اتجاه اللف المغزلي.

المركبة المقاسة وقيمتها المتوقعة شيئان مختلفان. فالقيمة المتوقعة متوسط على مجموعة إحصائية. حضّر عددًا كبيرًا من الإلكترونات في الحالة نفسها تمامًا، وقِس \(L_z\) مرة واحدة لكل إلكترون، ثم اجمع النتائج واقسمها على عدد التجارب. وكل نتيجة فردية تبقى إما \(+\hbar/2\) وإما \(-\hbar/2\)، لكن المتوسط قد يأخذ أي قيمة بينهما:
\begin{equation}
\langle L_z\rangle
=
\frac{\hbar}{2}\bigl(|\alpha|^2-|\beta|^2\bigr).
\end{equation}

إن تساوي سعتي حالتي الأعلى والأسفل يعطي لفًا مغزليًا موجَّهًا على طول أحد اتجاهي محور \(x\). والممثلان المطبَّعان هما\footnote{تحدد نسب المعاملات الاتجاهات. وتوفر العوامل \(1/\sqrt2\) التطبيع الذي يقتضيه \( |\alpha|^2+|\beta|^2=1\). وتعطي الأطوار النسبية \(+i\) و\(-i\) اتجاهين متعاكسين على محور \(y\)؛ أما إسناد أيٍّ من الإشارتين إلى \(y\) الموجب فيعتمد على اصطلاح المحاور.}
\begin{equation}
\lvert x;+\rangle
=
\frac{\lvert+\rangle+\lvert-\rangle}{\sqrt2},
\qquad
\lvert x;-\rangle
=
\frac{\lvert+\rangle-\lvert-\rangle}{\sqrt2}.
\end{equation}
وبالنسبة إلى أيٍّ من الحالتين، فإن القياس على طول \(z\) يعطي النتيجتين الممكنتين باحتمالين متساويين، ولذلك
\begin{equation}
\langle L_z\rangle
=
\frac12\left(\frac{\hbar}{2}\right)
+
\frac12\left(-\frac{\hbar}{2}\right)
=0.
\end{equation}
وتعطي الإشارتان النسبيتان الموجبة والسالبة اتجاهين متعاكسين على محور \(x\). وبالمثل،
\begin{equation}
\frac{\lvert+\rangle+i\lvert-\rangle}{\sqrt2},
\qquad
\frac{\lvert+\rangle-i\lvert-\rangle}{\sqrt2}
\end{equation}
تعطيان الاتجاهين المتعاكسين على محور \(y\). وتصف القيم المطبَّعة العامة لـ \(\alpha\) و\(\beta\) محاور اعتباطية. وإن قولنا إن اللف المغزلي يشير على طول محور معيّن يعني أن قياس المركبة على طول ذلك المحور يعطي نتيجة محددة. وبالنسبة إلى محور مختلف، تكون الحالة نفسها عمومًا تراكبًا.

تبقى نتائج القياس الفردية متكمّمة. وإنما متوسطاتها على المجموعات الإحصائية هي التي تتغير على نحو مستمر وتعيد إنتاج الهندسة المألوفة لمتجه كلاسيكي.

\section{اللف المغزلي، والإحصاء، والجسيمات المركبة}

تُعَدّ الكتلة والشحنة الكهربائية واللف الذاتي من بين الخصائص الرئيسية المستخدمة في توصيف الجسيم. ويرتبط اللف الذاتي أيضًا بالإحصاء. وعلاقة اللف الذاتي بالإحصاء مبرهنة في ميكانيكا الكم النسبية، مع أننا لن نبرهنها هنا:
\begin{equation}
\begin{aligned}
\text{half-integer spin} &\quad\Longleftrightarrow\quad \text{fermion},\\
\text{integer spin} &\quad\Longleftrightarrow\quad \text{boson}.
\end{aligned}
\end{equation}
لا يمكن لفرميونين متماثلين أن يشغلا الحالة الكمومية الكاملة نفسها. أما البوزونات المتماثلة فيجوز لها أن تشغل الحالة نفسها، ويمكن لأعداد كبيرة منها في حالة واحدة أن تسلك جماعيًا مثل موجة كلاسيكية. لذلك فإن الجسيمات ذات اللف الذاتي \(\tfrac12,\tfrac32,\tfrac52,\ldots\) هي فرميونات، في حين أن الجسيمات ذات اللف الذاتي \(0,1,2,\ldots\) هي بوزونات.

تمتد هذه القاعدة مباشرة إلى الجسيمات المركبة. إن جمع عزمين زاويين من نوع نصف صحيح يعطي عزمًا زاويًا صحيحًا. وعلى نحو أعم، فإن عددًا زوجيًا من المكوّنات الفرميونية يصنع مركبًا بوزونيًا، في حين أن عددًا فرديًا يصنع مركبًا فرميونيًا. وإضافة مكوّنات بوزونية لا تغيّر هذه القاعدة المتعلقة بالزوجية والفردية.

يحتوي الهيدروجين العادي على بروتون ذي لف ذاتي \(\tfrac12\) وإلكترون ذي لف ذاتي \(\tfrac12\). ولذلك فهو يحتوي على عدد زوجي من الفرميونات ويكون بوزونيًا ككل. ويحتوي الديوترون على بروتون ونيوترون، وكلاهما فرميون، وهو أيضًا بوزون. وتضيف ذرة الديوتيريوم إلكترونًا إلى الديوترون، فتكون هناك ثلاثة مكوّنات فرميونية على هذا المستوى من الوصف، ولذلك فهي فرميون.

للجسيمات والجسيمات المضادة الكتلة نفسها واللف الذاتي نفسه، كما أن شحناتها الكهربائية متعاكسة. البوزيترونيوم هو إلكترون وبوزيترون مرتبطان أحدهما بالآخر. وهو يحتوي على فرميونين، ولذلك فهو بوزوني.

\begin{classroomqa}
\lecturetimestamp{01:33:08}
\audiencequestion{هل تمرّ البوزونات خلال الأشياء، وهل ترتبط الموصلية الفائقة بسلوك أزواج الإلكترونات كبوزونات؟}
\lecturerresponse{الخاصية ذات الصلة ليست مجرد القدرة على المرور خلال المادة؛ بل هي قدرة البوزونات على شغل الحالة نفسها. ويقدّم الهيليوم فائق السيولة المثال الأبسط. تحتوي ذرة الهيليوم-\(4\) العادية على بروتونين ونيوترونين وإلكترونين، أي عدد زوجي من المكوّنات الفرميونية، لذا فالذرة بوزون. ويمكن لعدد كبير من ذرات الهيليوم أن يشغل حالة كمومية واحدة وأن يتحرك جماعيًا. وعندئذٍ يمكن تصوير المائع الفائق، على نحو تقريبي، على أنه موجة كلاسيكية تتكوّن من وضع عدد كبير من الذرات في الحالة نفسها. وهذه الصورة ليست دقيقة تمامًا، لكنها تلتقط الفكرة المطلوبة هنا. فالإلكترون في فلز هو فرميون، ولا يمكنه أن يشارك إلكترونًا آخر في حالة كاملة. أما أزواج الإلكترونات فهي مركبات بوزونية، وهنا يكمن ارتباط ذلك بالموصلية الفائقة. وهذا يثير مفارقة: إذا كانت الإلكترونات المكوِّنة لا تستطيع شغل الحالة نفسها، فكيف يمكن لعدة أزواج مركبة أن تشغل حالة واحدة؟ ويتطلب حل هذه المفارقة عرضًا أكثر تدقيقًا للحالات المركبة، ولذلك يُؤجَّل.}
\end{classroomqa}

\begin{classroomqa}
\lecturetimestamp{01:35:46}
\audiencequestion{بما أن البروتونات والنيوترونات مركبة، فلماذا لها لف ذاتي نصف صحيح؟}
\lecturerresponse{كلٌّ منهما يتكوّن من عدد فردي من الكواركات ذات اللف الذاتي \(\tfrac12\). وثلاثة مكوّنات فرميونية تصنع مركبًا فرميونيًا، ولذلك يمكن أن يكون للبروتونات والنيوترونات لف ذاتي نصف صحيح.}
\end{classroomqa}

\section{مبدأ استبعاد باولي والمدارات الذرية}

يوفّر الجدول الدوري الدافع المباشر لمبدأ استبعاد باولي. وليست القاعدة أن إلكترونين لا يمكنهما أبدًا أن تكون لهما الحركة المدارية نفسها. بل إن القاعدة هي أن إلكترونين لا يمكنهما شغل الحالة الكمومية الكاملة نفسها.

ولهذا التمييز أهمية حتى في أبسط الذرات. انطلاقًا من الهيدروجين، زِد الشحنة النووية وأضف إلكترونًا ثانيًا. يمكن للإلكترونين أن يشغلا المدار المكاني نفسه، ولكن لا بد أن تختلف حالتا لفّهما الذاتي. فقد يكون أحدهما ذا لف ذاتي إلى أعلى والآخر ذا لف ذاتي إلى أسفل بالنسبة إلى المحور المختار. وعندما يشغل إلكترونان مدارًا واحدًا، تتراكب حالتا لفهما الذاتي لتكوين الحالة المفردة للَّف الذاتي، بحيث ينعدم لفهما الذاتي الكلي. ولا يجوز لأكثر من إلكترونين شغل مدار مكاني واحد، لأن حالتي لف ذاتي مستقلتين فقط متاحتان.

\begin{classroomqa}
\lecturetimestamp{01:37:53}
\audiencequestion{هل يمكن تصوير اللفين الذاتيَّين المتعاكسين لإلكترونين في مدار واحد على أنهما تروسان متعاشقان؟}
\lecturerresponse{بوصفه تصويرًا تقريبيًا، نعم: فالتروس المتعاشقة تدور في اتجاهين متعاكسين. غير أن هذا مجرد تشبيه لاتجاهي اللف الذاتي المتعاكسين، وليس نموذجًا ميكانيكيًا حرفيًا للإلكترونين.}
\end{classroomqa}

يمكن استخدام كلمة ``مدار'' على مستويين مختلفين. فحالة الحركة المدارية، بالمعنى الواسع، هي الحالة التي نحصل عليها من حل معادلة شرودنغر مع إهمال اللف الذاتي. وهي تتضمن معلومات أكثر من مجرد تحديد الغلاف الذي يشغله الإلكترون.

\begin{classroomqa}
\lecturetimestamp{01:38:09}
\audiencequestion{هل يعني المدار أكثر من مجرد الغلاف الذي يشغله الإلكترون؟}
\lecturerresponse{نعم. فالحركة المدارية هنا تعني حالة شرودنغر للإلكترون عند إهمال اللف الذاتي. وهي تتضمن عدد الكم الرئيسي \(n\)، الذي يحمل معلومات شعاعية ويحدد تقريبًا بُعد الإلكترون عن النواة، وكذلك أعداد الكم للعزم الزاوي المداري. وفي الترميز الطيفي الذي يستخدمه الكيميائيون، تقابل \(S\) و\(P\) و\(D\) و\(F\) على الترتيب قيم العزم الزاوي المداري \(l=0,1,2,3\).}
\end{classroomqa}

وهكذا يصف عدد الكم الرئيسي والعزم الزاوي المداري جانبين مختلفين من الحركة. وتشير العلامات \(S,P,D,F\) إلى قيمة \(l\)، بينما يوفّر \(n\) معلومات شعاعية منفصلة. وقاعدة باولي هي أنه لا يجوز لإلكترونين أن تتطابق جميع أعدادهما الكمومية، بما في ذلك عدد الكم للَّف الذاتي.

\begin{classroomqa}
\lecturetimestamp{01:40:42}
\audiencequestion{هل يمكن لأعداد الكم أن تحدد عدد الإلكترونات التي قد تشغل كل غلاف في الذرة؟}
\lecturerresponse{نعم. فهي تحدد العدد الأقصى الممكن، لكنها لا تحدد عدد الإلكترونات التي تحتويها ذرة معينة فعلًا في ذلك الغلاف. ولن نتابع هنا المسألة العامة لعدّ الأغلفة.}
\end{classroomqa}

بالنسبة إلى عزم زاوي مداري ثابت \(l\)، فإن لعدد الكم للمركبة
\begin{equation}
2l+1
\end{equation}
من القيم الممكنة. وهذا هو العدّ نفسه للعزم الزاوي الذي استُخدم من قبل للَّف الذاتي. فحالة \(S\) لها \(l=0\)، ولذلك لها حالة مدارية واحدة. وحالة \(P\) لها \(l=1\)، ولذلك لها ثلاث حالات مدارية. وحالة \(D\) لها \(l=2\)، ولذلك لها خمس حالات مدارية.

ويضاعف اللف الذاتي كل عددٍ من هذه الأعداد المدارية، لأن كل مدار مكاني يمكن أن يحمل إلكترونًا واحدًا في كلٍّ من حالتي اللف الذاتي المتعاكستين.

\begin{classroomqa}
\lecturetimestamp{01:41:54}
\audiencequestion{هل تتسع حالة \(S\) لإلكترونين، وهل تتسع الحالات الثلاث \(P\) لستة أخرى؟}
\lecturerresponse{نعم. فالحالة المدارية الوحيدة \(S\) تتسع في الحد الأقصى لإلكترونين. والحالات المدارية الثلاث \(P\) تتسع لاثنين لكل منها، فتعطي ستة أخرى. لذا فإن حالات \(S\) و\(P\) قيد المناقشة تستوعب في المجموع \(2+6=8\) إلكترونًا.}
\end{classroomqa}

وتظهر البنية نفسها للعزم الزاوي التي تنظّم المدارات الذرية من جديد، مع بعض الاختلافات، عندما تتحرك الكواركات في منظومات مقيدة ويدور بعضها حول بعض.